\appendix
\iftoggle{icml}{\onecolumn}{}

\makeatletter
\@ifpackageloaded{lineno}{%
  \setlength{\linenumbersep}{10pt}
  \linenumbers
  \setlength{\linenumberwidth}{2.5em} % room for 3 digits
}{}
\makeatother

\mtcaddpart[Supplementary Materials]
\thispagestyle{empty}
{%
\bfseries\Large\vspace{-2.0em}%
\begin{center}
{%\toptitlebar{%
\ourtitleshort{}\\[0.25em]%
Supplementary Materials\\[0.25em]%
}%\bottomtitlebar
\end{center}%
\vspace{-1.0em}
}
\renewcommand\ptctitle{}
\mtcsetfeature{parttoc}{open}{}%\vspace{1.0em}\hrule}
\setlength{\ptcindent}{0pt}
\noptcrule
\parttoc[c]
\clearpage

\section{Notation}
\label{Appendix::Notation}
We list recurring notation in \cref{Table::Notation}. We use $i$ to index examples and $j$ or $k$ to index concepts.

\begin{table}[h]
    \centering
    \resizebox{0.75\textwidth}{!}{
    \begin{tabular}{ll}
        \toprule
        \textbf{Symbol} & \textbf{Meaning} \\
        \midrule
        $n$ & Number of examples in dataset. \\
        $m$ & Number of concepts. \\
        $\features_i \in \X \subseteq \R^d$ & Feature vector for example $i$ (e.g., pixels). \\
        $\concepts_i \in \C = \{0,1\}^m$ & Ground-truth concept vector for example $i$. \\
        $\outcome_i \in \Y = \{0,1\}$ & Ground-truth label for example $i$. \\
        $\detector_k: \X \to [0,1]$ & Concept detector for concept $k$; predicts $\Pr{c_{i,k} = 1 \mid \features_i}$. \\
        $\frontend: \C \to [0,1]$ & Front-end model mapping concepts to soft labels. \\
        $\pipeline: \X \to \Y$ & Full CBM pipeline outputting hard labels. \\
        $\baseline: \X \to \Y$ & Baseline model used for comparison (e.g., a black-box \DNN). \\
        $\outcomepred_i = \pipeline(\features_i)$ & Hard label prediction for example $i$. \\
        $\threshold \in [0, 0.5]$ & Abstention threshold for selective classification. \\
        $\abstain$ & Abstention symbol (model refuses to predict). \\
        \midrule
        \multicolumn{2}{l}{\textbf{Interventions and Metrics}} \\
        \midrule
        $t: \C \times \X \to \C$ & Intervention mapping modifying concept values given input $\features$. \\
        $S(\features) \subseteq \{1,\ldots,m\}$ & Subset of concepts selected for intervention for input $\features$. \\
        $|S(\features)|$ & Number of concepts requiring intervention for input $\features$. \\
        $\gain(\pipeline, \baseline)$ & Accuracy gain of CBM $\pipeline$ over baseline $\baseline$ after interventions. \\
        $\textsf{SelectiveAccuracy}(\pipeline)$ & Accuracy on non-abstained examples. \\
        $\textsf{Coverage}(\pipeline)$ & Proportion of examples where model does not abstain. \\
        $\wrk(\pipeline)$ & Net work automated (coverage minus intervention cost). \\
        $\incacc$ & Improvement in accuracy due to interventions (same model, before vs.\ after). \\

        \midrule
        \multicolumn{2}{l}{\textbf{Task-Specific Notation}} \\
        \midrule
        $\nBoard$ & Board size parameter for Sudoku (e.g., $\nBoard=3$ for $9 \times 9$ boards). \\
        $\crow{j}, \ccol{j}, \cblk{j}$ & Sudoku constraint concepts for row $j$, column $j$, block $j$. \\
        $\textcp{concept\_name}$ & Specific concept names (e.g., \textcp{FootShape}, \textcp{HasKnees}). \\
        $\outcome$ & True robot type in robots benchmark ($\textcp{G}$ for Glorp, $\textcp{D}$ for Drent). \\
        \bottomrule
    \end{tabular}
    }
    \caption{Notation used throughout the paper.}
    \label{Table::Notation}
\end{table}


%\clearpage
% \section{Dataset Noise Calculation Methodology}
% \label{Appendix::DatasetTable}
% We calculated concept noise statistic for a given concept $c$ as follows:

% \begin{itemize}
%     \item CUB-200: Counted the proportion of instances where $c$ was explicitly marked as uncertain/ambiguous (2 = guessing, and 3 = probably in the dataset).

%     \item AwA2: Counted the proportion of concept ratings $r\in[0,1]$ for concept $c$, later binarized, that were initially imperfect, i.e., $r\neq 0$ and $r\neq 1$.

%     \item CelebA: Noise rates from taken the study by \cite{lingenfelter2022quantitative} identifying annotation inconsistencies.

%     \item CheXpert: Counted the proportion of "uncertain" labels $u$ for concept $c$ being the result of automated extraction from medical reports, as marked in original dataset.

%     \item CEBaB: Counted the proportion of labels where there was no inter-rater agreement on the majority value of concept $c$
% \end{itemize}

% Then we computed:
% \begin{itemize}
%     \item Instance Noise = (data points with at least 1 noisy concept) / total data points
%     \item Concept Noise = average across all concepts of (noisy annotations for that concept) / (total annotations for that concept)
% \end{itemize}

% We couldn't get access to the OAI dataset and report only the values available in works using this dataset.

%% ======================================================================
%% B. Supplementary Material for Decision Support (Robots) → Section 3
%% ======================================================================
\clearpage
\section{Supplementary Material for Decision Support}
\label{Appendix::DecisionSupport}


\subsection{Dataset Parameters and Setup}
\label{Appendix::RobotsImageParams}
Here, we provide a detailed description of the properties of the \textds{robots} dataset.
\begin{itemize}
    \item \textfn{concepts}: head shape (square/round), body shape (square/round), knees (present/absent), antennae (present/absent), ears (square/triangle), mouth (closed/open), and foot shape with ten subtypes (flat trapezoid/flat rounded/flat square/flat five-sided/flat L-shaped/pointy trapezoid/pointy rounded/pointy square/pointy three-sided/pointy four-sided). Hands have six subtypes (three round and three angular) and feet have ten subtypes (five flat and five pointy). In the main concept set, we binarize hands (round vs. angular) and feet (flat vs. pointy). For the \textds{subconcepts} experiment, we replace the coarse \textcp{FootShape} concept with six \textcp{FootShape\_is\_Subtype} indicators corresponding to the six subtypes used in the skew constraints.
    \item \textfn{drop\_concepts}: elbows and hand shape are rendered but excluded from the concept bottleneck. Hand shapes have six subtypes (round circle/round oval/round oval variant/edgy triangle/edgy square/edgy trapezoid).
    \item \textfn{skew}: foot-shape subtypes are imbalanced with two dominant types (pointy four-sided and flat five-sided at $0.49$ each) and four rare types (pointy square/pointy rounded/flat square/flat trapezoid at $0.005$ each). The remaining four subtypes do not appear in the generated data.
    \item \textfn{drawing\_parameters}: two-tone color images with medium resolution (32px generator resolution).
    \item \textfn{samples\_per\_instance}: four color renderings per robot; color is treated as a spurious attribute.
    \item \textfn{dataset\_sizes}: train = 3{,}800, test = 10{,}000, with the remainder used for validation.
    \item \textfn{split\_procedure}: We draw a test set of 10{,}000 robot instances uniformly from the rendered-instance pool (unique robots $\times$ color renderings). From the remaining instances, we construct a training set of size 3{,}800 with minimum shares for six foot subtypes. Two dominant subtypes (one pointy, one flat) each have a 49\% minimum share, and four additional subtypes each have a 0.5\% minimum share. The remaining subtypes are unconstrained. Any leftover slots are filled uniformly from the remaining instances. The validation set is drawn from what remains after training selection.
    \item \textfn{label\_model}: We use the stochastic labeling model in \cref{Eq::RobotGT} with the weights, intercept, and scalar listed in \cref{Appendix::RobotsSetupParams}.
\end{itemize}

{\small
  \begin{align}
      \prob{\outcome=\textcp{G}\mid \concepts}
      = \sigma \Bigl(4.2\bigl(5\cdot\indic{\textcp{MouthType} = \textcpp{Closed}} +
  8\cdot\indic{\textcp{FootShape} = \textcpp{Pointy}} - 5\cdot\indic{\textcp{HasKnees}
  = \textcpp{True}} + 2\bigr)\Bigr)
      \label{Eq::RobotGT}
  \end{align}
  }

\subsection{Benchmark Instantiation Parameters}
\label{Appendix::RobotsSetupParams}
We list the parameter values used to instantiate the \textds{robots} benchmark reported in the main text.
\begin{itemize}
    \item \textfn{model\_type}: stochastic.
    \item \textfn{model\_features}: $\{\textcp{MouthType}: \textcpp{Closed},\ \textcp{FootShape}: \textcpp{Pointy},\ \textcp{HasKnees}: \textcpp{True}\}$.
    \item \textfn{model\_weights}: $\{\textcp{MouthType}: 5,\ \textcp{FootShape}: 8,\ \textcp{HasKnees}: -5\}$.
    \item \textfn{model\_intercept}: $2.0$. \quad \textfn{model\_scalar}: $4.2$.
    \item \textfn{concepts}:
    \textcp{HeadShape} $\in \{\text{square}, \text{round}\}$; \\
    \textcp{BodyShape} $\in \{\text{square}, \text{round}\}$; \\
    \textcp{HasKnees} $\in \{\text{false}, \text{true}\}$; \\
    \textcp{HasElbows} $\in \{\text{false}, \text{true}\}$; \\
    \textcp{HasAntennae} $\in \{\text{false}, \text{true}\}$; \\
    \textcp{EarShape} $\in \{\text{square}, \text{triangle}\}$; \\
    \textcp{MouthType} $\in \{\text{closed}, \text{open}\}$; \\
    \textcp{HandShape} $\in \{\text{round circle}, \text{round oval}, \text{round oval variant}, \text{edgy triangle}, \text{edgy square}, \text{edgy trapezoid}\}$; \\
    \textcp{FootShape} $\in \{\text{flat trapezoid}, \text{flat rounded}, \text{flat square}, \text{flat five-sided}, \text{flat L-shaped}, \text{pointy trapezoid}, \text{pointy rounded}, \\ \text{pointy square}, \text{pointy three-sided}, \text{pointy four-sided}\}$.

    \item \textfn{drop\_concepts}: \textcp{HasElbows}, \textcp{HandShape}, and fine-grained foot-shape subtypes. After dropping, the CBM receives 7 binary concepts (with \textcp{FootShape} binarized to flat vs.\ pointy).
    \item \textfn{skew}: foot-shape subtype probabilities are $0.005$ (pointy square), $0.005$ (pointy rounded), $0.49$ (pointy four-sided), $0.005$ (flat square), $0.005$ (flat trapezoid), and $0.49$ (flat five-sided). These six subtypes exhaust the probability mass; the remaining four subtypes do not appear in the generated data.
    \item \textfn{drawing\_parameters}: \textfn{image\_size} = medium (resolution 32), \textfn{color\_mode} = color.
    \item \textfn{samples\_per\_instance} = 4.
\end{itemize}

\subsection{Text Modality}
We also represent robots as text by prompting an LLM with the robot's attribute JSON and asking for a short, natural description (1–3 sentences), using the prompts in \cref{Appendix::DecisionSupport}. For each robot, we generate multiple text variants for the same underlying robot instance. For consistency across runs, we cached a set of LLM generations and reused them as fixed templates.

\paragraph{Concept noise via masking/omission}
To simulate ambiguity, we mask or omit subsets of attributes in the JSON and instruct the LLM to describe masked attributes only generically and to omit omitted attributes entirely. The remaining, unmasked attributes are described normally, so each caption encodes the same robot while selectively obscuring specific concepts.


\paragraph{Unstructured description prompt}
\begin{center}
\begin{tcolorbox}[promptbox, listing only, listing options={basicstyle=\ttfamily\small,breaklines=true}]
System prompt: You are concise and concrete. Use everyday language. Do not invent
locations or scenarios. \\

User prompt: Using the provided attributes, write a description in plain, spoken
language that sounds like a person describing an image they saw. Keep it
concise and natural, between 1 and 3 sentences. Avoid list-like phrasing,
avoid locations/situations, and focus only on what the attributes imply. \\

Attributes (JSON): \{...\}
\end{tcolorbox}
\end{center}

\paragraph{Masked/omitted prompt for concept noise}
\begin{center}
\begin{tcolorbox}[promptbox, listing only, listing options={basicstyle=\ttfamily\small,breaklines=true}]
System prompt: You describe simple robots from attribute lists. \\

User prompt: You are given attributes of a robot as JSON.
Write a single English sentence describing the robot.
Use natural language, not attribute names.
For attributes in 'masked', describe them only generically without revealing
their exact value.
For attributes in 'omitted', do not include any information about them.
All other attributes should be reflected in the sentence.
Do not add extra attributes that are not implied by the JSON. \\

JSON: \{...\}
\end{tcolorbox}
\end{center}

We provide a juxtaposition of different ways the data points can be represented in~\cref{Fig::Appendix-robots-modalities}.
% Dataset table
\begin{table}[H]
\centering
\label{tab:annotation_examples}

\resizebox{\textwidth}{!}{%

\begin{tabular}{p{0.75in} p{0.75in} p{0.75in} p{0.75in} p{0.75in} p{0.75in} c c}
\toprule

\multicolumn{3}{c}{\textbf{Ordered}} &
\multicolumn{3}{c}{\textbf{Shuffled}} &
\multicolumn{2}{c}{\textbf{Example Images}} \\

\cmidrule(r){1-3} \cmidrule(l){4-6}
\textbf{1 word} & \textbf{multiple words} & \textbf{Full phrases} &
\textbf{1 word} & \textbf{multiple words} & \textbf{Full phrases} & & \\
\midrule

``Head: square; Body: square; Color: teal--orange.'' &
``Head is square; Body is square; Color is teal--orange.'' &
``Body is square and head is matching; teal--orange color.'' &
``Color: teal--orange; Head: square; Body: square.'' &
``Color is teal--orange; Head is square; Body is square.'' &
``Teal--orange color; Head is square and body is matching.'' &

\includegraphics[width=1in, valign=m]{figures/robot_3836.png} &
\includegraphics[width=1in, valign=m]{figures/robot_3836_lowres.png} \\
\bottomrule
\end{tabular}%
}
\caption{Examples of different annotation styles and corresponding images.}
\label{Fig::Appendix-robots-modalities}
\end{table}


%% ======================================================================
%% C. Supplementary Material for Automation (Sudoku) → Section 4
%% ======================================================================
\clearpage
\section{Automation Details}
\label{Appendix::Automation}
\subsection{MIP Data Generation}
The Sudoku generator solves the following MIP for valid boards. The block size is $\nBoard{} \in \mathbb{N}$. Row, column, and value indices range over $\{1,\dots,\nBoard{}^2\}$. Random objective coefficients $w_{i,j,k} \sim \mathcal{U}(0,1)$ break ties across feasible boards.
%
\begin{samepage}
\begin{subequations}
\label{MIP::Sudoku}
\footnotesize
\renewcommand{\arraystretch}{1.5}
\begin{equationarray}{@{}c@{}r@{\,}c@{\,}l>{\,}l>{\,}l@{\;}}
%
\max_{x_{ijk}}\qquad & \sum_{i=1}^{\nBoard{}^2}\sum_{j=1}^{\nBoard{}^2}\sum_{k=1}^{\nBoard{}^2} w_{ijk}\, x_{ijk} \nonumber \\[1.5em]
\st
%
& \sum_{k=1}^{\nBoard{}^2} x_{ijk} & = & 1 & \textrm{for } i, j \in \intrange{\nBoard{}^2} & \mipwhat{one value per cell} \label{Con::Sudoku::Cell} \\
%
& \sum_{j=1}^{\nBoard{}^2} x_{ijk} & = & 1 & \textrm{for } i, k \in \intrange{\nBoard{}^2} & \mipwhat{each value once per row} \label{Con::Sudoku::Row} \\
%
& \sum_{i=1}^{\nBoard{}^2} x_{ijk} & = & 1 & \textrm{for } j, k \in \intrange{\nBoard{}^2} & \mipwhat{each value once per column} \label{Con::Sudoku::Col} \\
%
& \sum_{(i,j) \in \block_{p,q}} x_{ijk} & = & 1 & \textrm{for } p, q \in \intrange{\nBoard{}},\, k \in \intrange{\nBoard{}^2} & \mipwhat{each value once per $\nBoard{}\!\times\!\nBoard{}$ block} \label{Con::Sudoku::Block} \\
%
& x_{ijk} & \in & \{0,1\} & \textrm{for } i, j, k \in \intrange{\nBoard{}^2} & \mipwhat{set to 1 if cell in row $i$ and column $j$ has value $k$} \label{Con::Sudoku::Binary}
%
\end{equationarray}
\end{subequations}
\end{samepage}
%
where $\block_{p,q} = \{(i,j) \mid i \in \{(p{-}1)\nBoard{}+1,\ldots, p\nBoard{}\},\, j \in \{(q{-}1)\nBoard{}+1,\ldots, q\nBoard{}\}\}$ denotes the set of cells in block $(p,q)$.

Invalid boards start from a valid board and apply between $1$ and \textfn{max\_corrupt} corruption actions. These actions swap entries within selected rows or columns, swap whole rows or columns across bands or stacks, or duplicate a value within one row or column. We then recompute the row, column, and block concepts from the corrupted board.

\subsection{Dataset Parameters}
The \textds{sudoku} benchmark uses the following generation parameters.

\begin{itemize}
    \item $\nBoard{}$: block size. The default is $3$, which yields a $9 \times 9$ board with values in $\{1,\ldots,\nBoard{}^2\}$.
    \item \textfn{valid\_ratio}: fraction of valid boards. The default is $0.5$.
    \item \textfn{max\_corrupt}: maximum number of corruption actions used to make an invalid board. The default is $9$.
    \item \textfn{n\_samples}: number of boards to generate. The default is $1000$.
    \item \textfn{target\_accuracy}: minimum accuracy required on kept predictions for selective classification. The default is $0.9$.
    \item \textfn{data\_type}: representation type, either \texttt{tabular} or \texttt{image}. This determines the CBM embedding model.
    \item \textfn{handwriting}: for image datasets, handwriting variation is generated by varying font, font size, and font color.
    \begin{itemize}
        \item \textfn{starter\_percent}: fraction of cells shown as starter values. These use the default non-handwriting font and black color.
        \item \textfn{candidates}: candidate-value rendering settings, including the fraction of cells with candidates, the maximum number of candidates per cell, and the candidate color.
    \end{itemize}
\end{itemize}

\subsection{Benchmark Instantiation Parameters}
The reported \textds{sudoku} runs use the following instantiation.
\begin{itemize}
    \item $\nBoard{} = 3$ (board size $9 \times 9$).
    \item \textfn{valid\_ratio} = 0.5.
    \item \textfn{max\_corrupt} = 9.
    \item \textfn{data\_type} = \texttt{image}.
    \item \textfn{handwriting} = \texttt{true}.
    \item \textfn{candidates} = 0.22.
    \item \textfn{starter\_percent} = 0.35.
    \item \textfn{n\_samples} = 1000.
    \item \textfn{split} = 600/200/200 (train/val/test).
    \item \textfn{target\_accuracy} = 0.95.
\end{itemize}

%% ======================================================================
%% D. Supplementary Material for Demonstrations → Section 5
%% ======================================================================
\clearpage
\section{Supplementary Material for Demonstrations}
\label{Appendix::Demonstrations}

\subsection{Concept Leakage Metrics}
\label{Appendix::Leakage}

We compute two published leakage measures on our benchmark for all four architectures trained on \textds{human\_concepts} under the \regimePerfect{} regime: Concept-Task Leakage (\textsf{CTL}) and Interconcept Leakage (\textsf{ICL}) from~\citet{parisini2025leakage}.

\begin{table}[h]
\centering
\small
\begin{tabular}{lcc}
\toprule
Model & \textsf{CTL} & \textsf{ICL} \\
\midrule
\CBM{}     & 0.0523 & 0.0124 \\
\CEM{}     & 0.0982 & 0.0240 \\
\ProbCBM{} & 0.1084 & 0.0414 \\
\ECBM{}    & 0.0017 & 0.0000 \\
\bottomrule
\end{tabular}
\caption{Concept leakage measures on our benchmark. \ProbCBM{} has the largest \textsf{CTL} and \textsf{ICL} scores, while \ECBM{} has near-zero scores. The gap separates leakage in \ProbCBM{}, whose concept outputs carry task signal, from collapse in \ECBM{}, whose bottleneck carries little task signal.}
\label{Table::Leakage}
\end{table}

\subsection{Interventions}
We measure the likelihood of prediction $i$ changing upon interventions on concepts $S_i$ using \emph{responsiveness} $\flipProb(S_i)$:
%
\begin{align*}
\flipProb(S_i) = \sum_{a \in \{0,1\}^{|S_i|}} p(\va\mid S_i) \cdot \mathbbm{1}[\outcomepred_i \neq \outcomepred_i^a]
\end{align*}
%
where:
\begin{itemize}
\item $p(\va\mid S_i) = \prod_{j \in S_i} \conceptpred_{i,j}^{a_j} (1-\conceptpred_{i,j})^{1-a_j}$ is the probability that human intervention would set values of concepts $S_i$ to $\va$
\item $\outcomepred_i^a$ is the prediction when concepts $S_i$ are set to $\va$.
\end{itemize}
%
We only intervene on instance $i$ if the maximum responsiveness exceeds threshold $\flipThreshold$, i.e., $\flipProb(S_i^*) > \flipThreshold$, and intervene on concepts $S_i^* = \arg\max_{|S_i| \leq k} \flipProb(S_i)$.

\subsection{Model Architectures}
For the \textds{robots} image baseline, we use a three-block DNN. Each block includes ReLU and 2\texttimes2 max-pooling, and uses 16, 32, and 64 channels, respectively. The backbone is followed by a 128-unit fully connected layer, dropout at 0.5, and a two-class head. For \textds{robots} concept detection, we use a DNN backbone with one linear head per concept. We apply pooling for large images. For \textds{sudoku} images, concept detection uses a ViT-Base encoder with patch size 16 and input resolution 224. The concept head is a two-layer MLP. The CBM front-end is logistic regression with an L2 penalty. We set C to 1.0 and cap training at 1000 iterations.


\paragraph{Vision Transformer concept detector}
Our main \textds{robots} experiments use a DNN concept detector. We also use a ViT-Base model for the \textds{subconcepts}. The ViT-Base model is pretrained on ImageNet-21k (\texttt{google/vit-base-patch16-224-in21k}). This model is fine-tuned with AdamW using a learning rate of $5\times10^{-5}$, batch size 16, and 3 epochs. It uses a ViT-Base encoder with patch size 16 and input resolution 224, followed by one linear head per concept.

\subsection{Concept Missingness Handling}
We implement missingness at the concept level and mask missing labels during training. Under MCAR, each concept entry is independently missing with probability $p$. Under MNAR, missingness depends on the concept value. By default, the missingness probability is $1.5p$ if the concept is present and $0.5p$ if it is absent, with probabilities clipped to $[0,1]$. Missing labels are ignored by the default loss using an observation mask, so only observed concept labels contribute to training. We do not impute missing concepts unless explicitly stated.

\subsection{Optimization and Training Hyperparameters}
Unless otherwise stated, concept detectors are trained with Adam using a learning rate of $10^{-3}$ and batch size 64. We train for up to 10 epochs and use early stopping with patience 5 based on mean validation F1. In our robots image experiments, we use batch size 32 and train concept detectors for up to 50 epochs with patience 10. The CBM front-end uses logistic regression on concept labels and is trained to convergence with a maximum of 1000 iterations. The image DNN baseline in our scripts is optimized with AdamW and learning rate $5\times 10^{-5}$ using batch size 16, with the number of epochs set by the run configuration. \todo{Need to be more specific.}


% We adapt label-free concept bottleneck models (LFCBM) \citep{oikarinen2023label} for our text setting. The original framework, designed for vision, generates concept signals by comparing CLIP image embeddings to CLIP text embeddings. Our adaptation modifies this approach for a text-only domain.

% \emph{Adaptations to text}
% Since both concepts and inputs are text, we simplify the multi-modal CLIP signal and complex projection layer in 2 ways:

% \begin{itemize}
%     \item \emph{Text-based Concept Signal:} We replace the CLIP (image-text) signal with a purely text-based cosine similarity. We use a sentence-transformer to encode the input texts (creating matrix $H$) and the concept terms (creating matrix $E$). The raw concept signal $Z_t$ is the cosine similarity matrix $\text{cosine}(H, E)$.

%     \item \emph{Ridge Regression Projection:} The original method learns a projection adapter $W_c$ by optimizing a custom similarity function \citep{oikarinen2023label}. We replace this optimization step and instead learn a simple linear projection $W$ via Ridge regression. This maps the input text embeddings $H$ to the target concept similarities $Z_t$.
% \end{itemize}

% We then apply a threshold to the resulting projected scores to obtain the final hard concept vectors.


% Successful interventions require well-defined, predictive concepts. However, annotating concepts through human elicitation or machine discovery can introduce ambiguity and subjectivity.

% We measure how concept annotations can affect the benefits we expect from confirmation. We consider four regimes:
% \begin{itemize}

%   \item \emph{Perfect Annotation}: Concepts are known and labeled correctly; confirmation interventions are 100\% accurate.

%   \item \emph{Expert Annotations}: Concepts are known and labeled correctly for training. At test, intervention attempts are only $80 \%$ accurate which models expert uncertainty.

%   \item \emph{Machine Annotation}: We adapt the methods from \citet{oikarinen2023label} to identify concepts and label them without humans. We provide the same concept set as the other regimes.

%   \item \emph{DNN}: A neural network trained directly on text, representing a standard non-interpretable baseline.

% \end{itemize}

% % We test three systems: a DNN, a CBM, and a Conceptual Safeguard. The DNN serves as a black-box baseline to compare against. The Conceptual Safeguard uses a selective gate to abstain at low confidence and route those cases to targeted concept edits, whereas the CBM restricts inference to the concept vector so that test-time edits directly modify the decoder's input (see Section \ref{Sec::Preliminaries}).
% %to explore how a CBM's value hinges on the quality and easy detectability of its underlying concepts.

% To simulate the noisy labels common in subjective tasks, we degraded the textual information for the key concepts in 70\% of the data. We achieve this by replacing a precise descriptor (e.g., "triangular feet") with an ambiguous phrase (e.g., "has feet"), making the concept impossible for the model to detect reliably from the text. This masking was done at random. To isolate the value of intervention, we assume an expert can view the image of the ground-truth robot and specify the true concept value. We limit the number of interventions per "expert" to 2 as the maximum reasonable level of intervention per instance.



% % \begin{wrapfigure}[14]{R}
% \begin{wrapfigure}{R}{0.38\linewidth}

%     \centering
%     \includegraphics[width=\linewidth]{figures/acc_vs_budget_swapGeneric_footOpenPointy_balanced_v2_combined_no_machine.pdf}
%     \caption{Intervention accuracy by edit budget under concept noise. CBMs can outperform baselines with intervention, but require supervision.}
%     \label{fig::InterventionBudgetNoise}
% \end{wrapfigure}


% In \cref{fig::InterventionBudgetNoise}, we see that CBMs can outperform standard DNNs with interventions.
% Initially, CBMs for \emph{Expert} and \emph{Perfect Annotations} and  start at 72\% accuracy, whereas the DNN outperforms it with 78\% accuracy. Both models are limited by the lack of signal in the main data, with much of the key signal masked. However, after making just one targeted concept intervention per sentence under the \emph{Perfect Annotation} regime, the order is clearly reversed. The CBM increases to 81\%, while the black-box baseline cannot be modified without retraining.

% \newcommand{\Gain}{\mathrm{Gain}}
% \newcommand{\acc}{\mathrm{acc}}
% \newcommand{\ConfirmationCost}{\mathrm{ConfirmationCost}}

% We see, however, that the benefit of the CBM system decreases with the reliability of the expert, as shown in \ref{fig::InterventionBudgetNoise}. Both regimes require human attention on more than half of all instances (63.9 \%) with an $\textrm{InterventionRate}(f, \verify)$ = 26\% of robot instances. This level of supervision yields a gain \(\Delta\text{acc}\) of less than five percent (+ 2.64\%). This relatively small gain in performance falls further when we simulate a more realistic level of expert accuracy, as checks inflate \(\ConfirmationCost\) without commensurate \(\Delta\text{acc}\). The $\Delta\text{acc}$ from a single intervention decreases immediately to 7.1\% (from 8.9\%) and tapers as the number of interventions increase. At $k = 2$ in the \emph{Expert} regime, the $\Delta\text{acc}$ from $k = 1$ falls to +$3.59\%$.

% The positive effect of interventions is further reduced when concepts are not clearly defined. For the \emph{Machine Annotation} regime, we adapt the image-based concept discovery method from \citep{oikarinen2023label} for our text-based task (see \cref{Appendix::Demonstrations} for details). The chosen concepts can have overlap with true concepts (e.g \textfn{barrel-shaped} = feet), but at the same time admit spurious features due to the heavy masking (e.g \textfn{curved} = frame). This leads to a prediction of almost entirely the majority class, with an accuracy of 72\%, but a \textrm{Balanced Error Rate} of 0.5. We do not attempt to simulate interventions here, as there is no true "ground-truth" concept we can recover automatically.

% One of the core arguments for the adoption of a CBM is its ability to be "intervenable". But as we show, this is not an automatic guarantee, and is fragile and contingent on the quality of the concept annotations themselves. When concepts are poorly defined or unreliably measured, interventions fail to provide meaningful gains, and the additional value of a CBM becomes unreliable.

\subsection{Label Noise and Domain Knowledge}

\begin{wrapfigure}[15]{R}{0.33\linewidth}
\centering
\vspace{-1em}
\includegraphics[
width=\linewidth%
]{figures/domain_acc_update.pdf}
\caption{Accuracy of CBMs (dots) and DNNs (squares) trained at varying label noise probabilities. The top line indicates performance when we manually specify the front-end model.}
\label{Fig::FrontEndModel}
\end{wrapfigure}
%\paragraph{Specifying front-end model}


Datasets for automation tasks can include label noise (e.g., due to human disagreement or error). However, the ground-truth mapping from concepts to labels is well known in such tasks. Although the learned model may be influenced by the noise and therefore inaccurate, we can leverage task-specific knowledge by manually specifying the front-end model of the CBM to mitigate this effect. To investigate this mitigation method, we utilize the same \texttt{tabular} variant setup of the \texttt{sudoku} dataset as in \cref{Sec::Sudoku} and add label noise to the training data. We adjust its prevalence through the noise parameter $\rho$ and measure the performance of the remaining test set. We show the held-out test set performance for models trained under $\rho = 0.05, 0.1, 0.15, 0.2, 0.25$ and compare it to the specified model from domain knowledge (i.e., front-end model: $y_i = \prod_j c_{i,j}$).

We see in~\cref{Fig::FrontEndModel} that the CBM with a specified model outperforms a CBM with a learned model, while both outperform a DNN. Since the CBM uses concept predictions to determine the final label, the noise from the training label will have no influence on the specified front-end model, hence the accuracy remains stable at 94.8\% at each value of $\rho$. The learned front-end model declines slightly as $\rho$ increases - e.g., 88.6\% when $\rho = 0.05$, and 87.5\% accuracy where $\rho = 0.3$. The concept predictions help mitigate the full effects, but the learned front-end model has still been affected. However, the DNN lacks the advantages of concept-level predictions, and all the weights directly absorb label noise, thus leading to a steep decline from 88.6\% at $\rho = 0$ to 69.25\% at $\rho = 0.3$.
